\documentclass[11pt]{article}

\usepackage[margin=1in]{geometry}
\usepackage{booktabs}
\usepackage{array}
\usepackage{amsmath}
\usepackage{hyperref}
\usepackage{xcolor}
\usepackage{longtable}
\usepackage{parskip}
\usepackage{caption}
\usepackage{fancyhdr}
\usepackage{titlesec}
\usepackage[htt]{hyphenat}

\titleformat{\section}{\normalfont\Large\bfseries}{\thesection}{1em}{}
\titleformat{\subsection}{\normalfont\large\bfseries}{\thesubsection}{1em}{}

\pagestyle{fancy}
\fancyhf{}
\fancyhead[L]{Prediction Accuracy Report v2}
\fancyhead[R]{\thepage}
\renewcommand{\headrulewidth}{0.4pt}

\hypersetup{
    colorlinks=true,
    linkcolor=blue,
    urlcolor=blue,
    citecolor=blue,
    pdftitle={Prediction Algorithm Accuracy Report v2 - Post-Remediation},
    pdfauthor={Sports Analytics Project}
}

\newcommand{\code}[1]{\allowbreak\texttt{\small #1}\allowbreak}

\title{\textbf{Prediction Algorithm Accuracy Report — v2} \\ \large Post-Remediation: What Was Tried, What Shipped, What's Next}
\author{Sports Analytics Project}
\date{September 2026}

\begin{document}

\maketitle

\begin{abstract}
\noindent This is a follow-up to the initial backtest report (\code{prediction-accuracy-report.pdf}, same
directory), which identified five recommendations for improving the platform's four prediction algorithms.
All five were acted on. Two produced real, shipped improvements to system trustworthiness (an ongoing
accuracy-monitoring script, and honest UI framing of the weaker margin prediction). One was a documentation
fix. The remaining two — re-tuning Elo's constants and adding an opponent-defense adjustment to the
player-point predictors — were properly tested against held-out data and \emph{did not generalize}: both
looked marginally better on the data used to select them, then performed no better (or slightly worse) on a
chronological validation split. Per this project's existing practice of not presenting overfit results as
reliable, neither was shipped. This report documents why, and — more usefully — identifies where the actual
remaining headroom is, so future effort targets higher-leverage changes instead of repeating these two.
\end{abstract}

\tableofcontents

\section{What Changed Since v1}

\begin{table}[h]
\centering
\small
\renewcommand{\arraystretch}{1.3}
\begin{tabular}{@{}>{\raggedright\arraybackslash}p{0.27\textwidth}>{\raggedright\arraybackslash}p{0.12\textwidth}>{\raggedright\arraybackslash}p{0.55\textwidth}@{}}
\toprule
Recommendation & Outcome & Detail \\
\midrule
\#1 Accuracy-tracking job & \textbf{Shipped} & \code{check\_accuracy.py} in \code{apps/predictor} — read-only, re-runnable, reuses \code{elo.py}/\code{four\_factors.py} directly \\
\#2 Opponent-adjustment for player predictors & \textbf{Not shipped} & Backtested; edge too small and didn't hold on a validation split (\S\ref{sec:opponent}) \\
\#3 Re-tune Elo constants & \textbf{Not shipped} & Grid-searched; train-split winner validated worse than the defaults (\S\ref{sec:elo-tuning}) \\
\#4 De-emphasize Four Factors margin in UI & \textbf{Shipped} & \code{GameDetailPage.tsx}, \code{PredictionsPage.tsx} — margin is now visually secondary with an explanatory tooltip \\
\#5 Fix stale ``seed data'' comments & \textbf{Shipped} & \code{four\_factors.py} and its README — now describe the live full-season state \\
\bottomrule
\end{tabular}
\caption{Disposition of all five v1 recommendations.}
\end{table}

Three of five were low-effort, unambiguous wins with no tradeoffs. The other two were genuine attempts at
improving raw model accuracy, tested with the same discipline as everything else in this codebase (walk-forward,
leak-free, validated on held-out data) — and came back negative. That is a normal and expected outcome in
applied modeling, not a sign anything was done wrong; \S\ref{sec:why-negative} explains specifically why each
one was unlikely to work \emph{before} even seeing the results, with the benefit of hindsight, and
\S\ref{sec:whats-next} gives concrete, higher-leverage alternatives.

\section{Backtested Accuracy: Unchanged}

Because neither of the two accuracy-focused experiments was shipped, and the three shipped changes were
documentation/monitoring/UI-only, \textbf{every backtested accuracy number from v1 is still exactly correct
today} — re-run in full against the live database while writing this report and confirmed identical:

\begin{table}[h]
\centering
\small
\begin{tabular}{lrr}
\toprule
Model & Metric & Value (unchanged from v1) \\
\midrule
Elo win probability & Brier score & 0.2146 \\
Elo win probability & Accuracy & 65.2\% (vs. 55.5\% naive) \\
Four Factors margin & MAE & 12.52 pts (vs. 13.36 naive) \\
Fantasy points & MAE & 7.74 (vs. 8.25 naive) \\
Scorer points & MAE & 4.45 (vs. 4.64 naive) \\
\bottomrule
\end{tabular}
\caption{All four models' backtested accuracy is unchanged from the v1 report — see that report for full
methodology, calibration tables, and permutation test results, all still current.}
\end{table}

\section{Why the Two Tuning Attempts Came Back Negative}
\label{sec:why-negative}

\subsection{Elo constant re-tuning}
\label{sec:elo-tuning}

\textbf{Method.} A chronological 70/30 train/validation split (first 924 games to grid-search
\code{HOME\_COURT\_ADVANTAGE\_ELO} and \code{K\_FACTOR} by Brier score, last 397 games held out purely to
check generalization — a random split would be invalid here since Elo's rating updates are inherently
sequential).

\textbf{Result.} The train-split grid search's winner (HCA=30, K=30) beat the literature-default (HCA=75,
K=20) on the training data (Brier 0.2233 vs.\ roughly 0.2251), but validated \emph{worse} on the held-out
split: 0.1839 vs.\ 0.1829 Brier. That is the textbook signature of overfitting — a configuration that looks
better only because it was chosen using the exact data it's being judged on.

\textbf{Why, in hindsight, this was unlikely to work.} The training-split Brier surface across the entire
tested HCA range (0 to 100) was quite flat — 0.2233 to 0.2251, a spread of only 0.0018. When a hyperparameter
surface is that flat, the ``optimal'' point found by a grid search is mostly noise: small sampling
fluctuations in which specific games fell in the training window are enough to shift where the minimum
lands. A single NBA season (1,321 games, effectively $\sim$44 games per team-pair-orientation once split by
home/away) is simply not a lot of independent evidence about a single global home-court constant — the
literature-derived defaults were themselves fit on vastly more games across many seasons, so it would be
mildly surprising if one season's worth of this project's data could improve on them.

\subsection{Opponent-defense adjustment}
\label{sec:opponent}

\textbf{Method.} The simplest possible opponent adjustment — scale a player's baseline recency-weighted
prediction by the ratio of the upcoming opponent's leak-free running points-allowed average to the
leaguewide average at that point in the season — backtested against the full season, then re-checked on a
chronological 30\% validation-only slice.

\textbf{Result.} Full-season MAE improved by 0.006 (fantasy points) and 0.002 (raw points) — both
negligible. On the validation-only slice, the fantasy-points edge shrank further and the raw-points edge
\emph{flipped slightly negative} ($-0.0001$).

\textbf{Why, in hindsight, this was unlikely to work.} Team-level points-allowed only spans roughly $\pm 9\%$
across the entire league (106.9 to 126.0 points allowed per game, league mean 115.3). A $\pm 9\%$
multiplicative nudge is small next to a single player's game-to-game scoring variance (RMSE $\approx$ 10
fantasy points / $\approx$ 6 raw points) — the adjustment was always going to be a small correction relative
to the noise it's competing with. This isn't evidence that opponent strength doesn't matter to basketball
outcomes; it's evidence that \emph{team-wide} points-allowed is too coarse a proxy for it at the
individual-player level. A stronger signal (e.g. points allowed to players at the same position specifically,
or a proper defensive rating that adjusts for pace) might carry more information, but that's a materially
different, larger effort than the one-line adjustment tested here — see \S\ref{sec:whats-next}.

\section{Where the Real Headroom Is}
\label{sec:whats-next}

Both negative results above share a pattern worth naming explicitly: \textbf{they were both attempts to
extract more signal from data the models already have}, just recombined differently (better-tuned constants,
a coarser opponent proxy). That data was already close to fully exploited by the existing simple models. The
remaining, larger opportunities are in \textbf{data the models don't have access to at all yet} — these are
not guaranteed to work either, but they target genuinely unused information, which is a fundamentally
different (and more promising) kind of bet than re-weighting existing inputs.

\begin{enumerate}
    \item \textbf{Minutes / role-change signal for the player-point predictors.} Neither
    \code{predict.py} nor \code{game-detail.service.ts}'s predictor has any feature for playing time — a
    player whose minutes just changed (returning from injury, moved into the starting lineup, a coach's
    rotation change) will be mispredicted by a pure box-score average until enough new games accumulate to
    catch up organically. This is a bigger lever than opponent adjustment because minutes changes are a much
    larger relative swing in a player's output than an opponent's defensive quality is. Concretely: track a
    player's own minutes-per-game trend (the same recency-weighting technique already used for points would
    work directly for minutes) and either use it as an additional regression feature or predict a
    per-minute rate and multiply by a minutes forecast.
    \item \textbf{Offensive/defensive rebound split for Four Factors.} This is the one item from the
    original report requiring a real (if modest) ingestion/schema change: \code{PlayerGameStat} currently
    tracks only total rebounds, not an offensive/defensive split, which is why the current model uses 3 of
    Oliver's 4 factors and explicitly documents the gap. Adding that split is the single most direct way to
    close some of the Four Factors margin model's remaining error, since it's a genuinely new dimension of
    information (not a recombination of what's already there, unlike the two experiments above).
    \item \textbf{More seasons of data for Elo.} \S\ref{sec:elo-tuning}'s finding was specifically that one
    season isn't enough evidence to beat a good literature prior — that is a data-volume problem, not a
    methodology problem. Re-attempting the same grid search once 2-3 seasons are loaded is likely to behave
    very differently (a much less flat, better-identified Brier surface), and is a ``wait and re-run,'' not
    a ``design something new,'' item.
    \item \textbf{Home/away split for the player-point predictors.} Simple to add (the data already exists
    on every \code{Game} row, same as opponent identity was), well-established in the wider basketball
    analytics literature, and untested by any of the work in this report or its predecessor — a reasonable
    next candidate to backtest, though given \S\ref{sec:opponent}'s result, it should be validated the same
    way (train/validation split, not just full-season MAE) before being treated as a win.
\end{enumerate}

\section{Process Note}

Everything in this report — including the two negative results — was produced using the same walk-forward,
leak-free, held-out-validation discipline the underlying prediction code itself already enforces (see
\code{four\_factors.py} and \code{games.e2e-spec.ts}'s existing data-leakage tests). Two specific ideas not
generalizing is a normal rate of return for this kind of iteration, not a signal that the approach is wrong
— the value was in finding out \emph{before} shipping either change, via the same discipline that caught a
real future-data leak in this codebase once before (documented in \code{four\_factors.py}'s own module
docstring). The three items in \S\ref{sec:whats-next} are concrete next bets, ranked by how much genuinely
new information they'd add rather than by how easy they'd be to implement.

\end{document}
